\chapter{QDOAS Overview} \label{cha:intro}
\section{Introduction}
The experience of the \gls{bira} in the development and improvement of algorithms for the retrieval of trace gas concentrations goes back to the early 1990s, with atmospheric research activities using ground-based UV-Visible spectrometers aiming at the long-term monitoring of minor components involved in the catalytic destruction of the ozone layer or in anthropogenic pollution.

WinDOAS, the first program developed at \gls{bira} in 1997, knew a success story due to a friendly user interface completed with some po\-wer\-ful \glslink{doas}{DOAS} tools.  This program, extensively validated through different campaigns, has been used worldwide and for many different DOAS applications (mainly for ground-based and satellite applications).

QDOAS is a cross-platform implementation of WinDOAS: the software is portable to Windows and Unix-based operating systems whereas WinDOAS was designed only for Windows. The user interface and the engine of QDOAS are similar to those of its predecessor.  WinDOAS is  no longer supported.

QDOAS has been developed in collaboration with \href{http://www.stcorp.nl/}{S[\&]T}.  The graphical user interface is built on the Open-Source version of the Qt toolkit, a cross-platform application framework, and Qwt libraries.

This document describes the QDOAS user interface and dialog boxes for configuring the software.  It completes the online help provided with the package.  The algorithms that are used in the software are also summarized.  The document assumes that users already have a minimum of experience with \gls{doas} applications and retrieval.

\section{Main QDOAS features}
\subsection{General features}
\begin{qdoasitem}
\item The main components of the \gls{gui} are organized in multi-page panels with a fixed arrangement and tab-switched access to the different pages;
\item The application is based on a tree structure;
\item All the spectral windows are processed in one shot and there is a tab-switched access between the different fitting windows;
\item Large amount of files can be processed in one shot;
\item Support different spectra file formats (see section \ref{sec:intro-formats}) and specificities for satellites, ground-based and airborne measurements are accounted for;
\item On line help in HTML format
\end{qdoasitem}
\subsection{Plot}
\begin{qdoasitem}
\item Visualization of spectra and the results in different tab pages;
\item Possibility to set plot contour and style;
\item Interactive plot mode (zooming, overlay of an existing ASCII file, possibility to fix the scaling of the plot, \ldots): activated by right clicking the title of the plot ;
\item Export of plot in different portable image formats (png, jpg)
\end{qdoasitem}
\subsection{Analysis}
\begin{qdoasitem}
\item DOAS/intensity fitting modes;
\item shift/stretch (in nm) fully configurable for any spectral item (cross section or spectrum);
\item automatic reference selection from the input spectra file taking account for specificities for satellites or MAXDOAS measurements (scans);
\item possibility to constrain the slant column density of a molecule to the value found in a previous window;
\item possibility to filter spectra and cross sections before analysis (supported filters include Kaiser, gaussian, boxcar, Savitsky Golay…);
\item possibility to define gaps or automatically remove spikes within fitting intervals (e.g. to eliminate bad pixels);
\item possibility to fit an instrumental offset;
\item possibility to define several configurations of spectral windows under a project;
\end{qdoasitem}
\subsection{Calibration And Slit Function Characterization}
\begin{qdoasitem}
\item wavelength calibration and instrumental slit function characterization using a \gls{nlls} fitting approach where measured intensities are fitted to a high resolution solar spectrum degraded to the resolution of the instrument. The fitting method (DOAS or intensity fitting) can be different from the method used in the analysis;
\item possibility to correct for atmospheric absorption and Ring effect;
\item supports different analytical line shapes, as described in section \ref{sec:wavel-calibr}, page \pageref{sec:wavel-calibr};
\item allows approximating the initial shift of the NLLS procedure by calculating the best correlation between the spectrum to calibrate and the solar spectrum (see \nameref{sec:projects_calibration_preshift}, page \pageref{sec:projects_calibration_preshift});
\end{qdoasitem}

\subsection{Cross sections handling}
\begin{qdoasitem}
\item possibility to calculate differential absorption cross sections (by orthogonalization or high-pass filtering);
\item possibility to multiply cross sections with wavelength dependent \glspl{amf};
\item possibility to fix the column density of any selected species;
\item possibility to convolve cross sections in real time using a user defined slit function or the information on calibration and slit function provided by the wavelength calibration procedure;
\item possibility to handle differences in resolutions between measured and control spectra;
\end{qdoasitem}

\subsection{Output}
The output is fully configurable. Analysis results and various data related to the measurements can be saved to ASCII or netCDF files.

\subsection{Tools}

The package includes five independent executables:

\begin{qdoasitem}
  \item qdoas: the main user interface;
  \item convolution: the convolution/filtering tool;
    \begin{itemize} \itemsep -2pt
    \item standard and \I0-corrected convolutions are supported;
    \item possibility to create an effective slit function taking into account the (finite) resolution of the source spectrum (using a FT deconvolution method);
    \item asymmetric line shapes, wavelength dependent slit functions;
    \end{itemize}

  \item ring: The ring tool calculates Ring effect cross sections (Rotational Raman Scattering approach);
  \item usamp: This tool generates undersampling cross sections;
  \item doas\_cl: a powerful command line tool that applies on qdoas, convolution, ring and usamp configuration files (XML format).
  \end{qdoasitem}
Convolution, ring and undersampling tools can be also called from the QDOAS user interface.

\section{Supported spectra file formats} \label{sec:intro-formats}
\subsection{Satellite Measurements}
Spectra measured by the following satellite instruments can be processed by QDOAS:

\subsubsection{GOME (ERS-2)}
QDOAS supports the GOME Level 1 Version 5.1 NetCDF format (\url{https://doi.org/10.5270/ER2-ua38y2m}), as well as
the legacy ASCII file format.  For users who wish to work with legacy
GOME data, we recommend to use a more practial binary file format
created by a modified version of the \glslink{gdp}{GDP} implemented at
\gls{bira} for the automatic selection of the reference spectrum.
Contact the authors for further information about this format.

\subsubsection{SCIAMACHY (ENVISAT)}
Routines to read the \glslink{sciamachy}{SCIAMACHY} PDS file format
have been kindly provided by IFE/IUP University of Bremen.  SCIAMACHY
Level 1C versions up to 7.04 (the SCIAL1C utility converts files from
L1 to L1C).

\subsubsection{GOME-2 (MetOp A/B/C)}
QDOAS uses the \glslink{coda}{CODA} library to read spectra from the
GOME-2A/B/C instruments.  In order to make this work, CODA needs a
.codadef file which describes the GOME-2 data format structure.
Download the file EPS-20250709.codadef from
\url{https://github.com/stcorp/codadef-eps/releases}, and make sure
the \verb|CODA_DEFINITION| environment variable contains the path
where the .codadef file is located (see
\url{https://stcorp.github.io/coda/doc/html} for details).

\subsubsection{OMI (AURA)}
QDOAS can work with \glslink{omi}{OMI} L1B collection 3 data, which
uses the \glslink{hdfeos}{HDF-EOS2} format, as well as the improved
collection 4 data, which is available in NetCDF format.

\subsubsection{TROPOMI (S5P)}
QDOAS can process TROPOMI Level 1B radiance
(\url{https://doi.org/10.5270/S5P-kb39wni}) and irradiance
(\url{https://doi.org/10.5270/S5P-mhtbru8}) products.

\subsubsection{OMPS (Suomi NPP/NOAA-20)}
QDOAS supports OMPS Nadir Mapper Level 1B data (\url{https://dx.doi.org/10.5067/DL081SQY7C89}).

\subsubsection{GEMS (GK-2B)}
QDOAS can process GEMS Level 1 radiance and irradiance data, which are
available upon request from \gls{nier}
(\url{https://nesc.nier.go.kr/en/html/index.do}).

\subsubsection{TEMPO (Intelsat 40e)}
QDOAS supports TEMPO Level 1 radiance
(\url{https://dx.doi.org/10.5067/IS-40e/TEMPO/RAD_L1.004}) and
irradiance (\url{https://dx.doi.org/10.5067/IS-40e/TEMPO/IRR_L1.004})
products, currently at ``provisional'' version 4.

\subsection{Ground-based Measurements}

The most popular file formats supported by QDOAS for ground-based measurements are:

\subsubsection{SAOZ}

Système d'Analyse par Observation Zénitale developed by the aeronomy lab of the CNRS (JP Pommereau, F Goutail,...), France, and largely used in the \glslink{ndacc}{NDACC} network for stratospheric total ozone and \NO2 monitoring.   Both PCD/NMOS 512 and EFM 1024 (SAM) formats are supported.

\subsubsection{MFC format, DOASIS} \label{sec:MFC_DOASIS}

Both binary and std (ASCII) formats are produced by the well-known \glslink{doasis}{DOASIS} program designed by the atmospheric research group at IUP Heidelberg, Germany and widely used in the DOAS community.

This format can also be produced by other programs or converted from other formats, e.g.\ the Ocean Optics SpectraSuite software. QDOAS requires the size of the detector.
The MFC binary format generated by DOASIS is supported by QDOAS since version 2.111.  The specificity of this format is that spectra are saved in individual files that are not easy to manage mainly for the automatic reference selection (zenith-sky spectrum measured at the minimum SZA of the day or zenith-sky of the scan for MAXDOAS measurements).
To browse and process them, it is recommended to insert folders instead of individual files.  Browsing records requires the use of the file selection combobox in the toolbar (cfr page \pageref{sec:interface_toolbar} ) as one file contains only one record in this format.  Nevertheless, the
automatic reference selection is automatically performed by browsing all the files of the current directory.  This could make the analysis heavier and a conversion to the BIRA-IASB MFC Binary format is recommended (using the \texttt{MFC\_Std2Bin} converter available on the FTP server with QDOAS)

\subsubsection{MFC binary, BIRA-IASB} \label{sec:MFC_bin}

Because spectra are saved in individual files with automatic numbering, MFC STD format is not very suitable to browse spectra and for efficient automatic reference selection (zenith of the scan for MAXDOAS data).  For this reason, BIRA-IASB has developed its own MFC binary format
in which all spectra of a same day including dark currents and offsets measured the night are saved in a unique file.  When loading a file, measurements are automatically corrected by the averaged dark currents and offsets.  The \texttt{MFC\_Std2Bin} utility that converts
MFC STD files in the MFC BIRA-IASB binary format is available with the QDOAS package since version 2.107.  Syntax to use from command line is

\begin{quote}
\begin{Verbatim}[fontfamily=courier,fontsize=\relsize{-2}]
MFC_Std2Bin <config file> <input file or path> <output path>

where :

  config file        : is the name of a file with some options
  input file or path : the name of the MFC file or path to process
  output path        : the path where to save the data

It is recommended to apply on dark current/offset files or paths
before spectra.  Daily output files are automatically generated
using the dates of measurement.  Yearly directories are created
in the output path.

The config file is an ASCII file including lines starting with the
following keywords :

   date format=...      : the format of the date, for example YYYYMMDD
   spectrum=...         : the string identifying spectra measurements
   offset=...           : the string identifying offset measurements
   dark=...             : the string identifying dark current measurements
   output prefix=...    : the prefix of output file name (the date in
                          YYYYMMDD format will follow)
   output extension=... : the file extension to complete the output file name

The program will use the strings identifying spectra, offset and dark
current to attribute a measurement type number to spectra :

   1 for off axis
   3 for zenith sky
   4 for dark currents
   8 for offsets

Example of config file :

   date format=MMDDYYYY
   spectrum=Measurement Spectrum
   offset=Offset Spectrum
   dark=Dark Current Spectrum
   output prefix=BX
   output extension=bin

All files of 07/01/2014 will be saved in <output path>/2014/BX_20140107.bin
\end{Verbatim}
\end{quote}

\subsubsection{MKZY PAK}

developed at the Chalmers University of Technology, Goteborg, Sweden (MANNE Kihlman and ZHANG Yan) and used in the NOVAC network.

\subsubsection{OCEAN OPTICS}

very simple ASCII file format developed by the spectrometers manufacturer.

\subsection{The ASCII File Format}
Other formats are specific to the different institutes that developed them.  If a format is not supported by QDOAS, it should be possible to convert the files to ASCII. When the ASCII format is selected, spectra can be provided in the file one record per line \textbf{(line format)} or one spectral value per line \textbf{(column format)} or several spectra in columns (from version 2.107). According to the checked flags in the \textbf{Read from file} group box, the following information are expected strictly in the given order~:
\begin{itemize} \itemsep -6pt
\item Solar Zenith Angle
\item Azimuth Viewing Angle
\item Elevation Viewing Angle
\item Date in the DD/MM/YYYY (day/month/ year) format
\item Fractional time
\end{itemize}

\textbf{In the column format, angles had to be given on the same line until version 2.106.
From version 2.107, the column format accepts also matrices of spectra and angles have to be given in separate lines.}  The file can contain
several matrices of spectra but the number of columns has to be the same for all data sets.

\begin{quote}
\begin{Verbatim}[fontfamily=courier,fontsize=\relsize{-2}]
Example of ASCII file in column format :

* data set with 4 individual spectra
* line 1 : solar zenith angle
* line 2 : elevation angle (azimuth viewing angle not provided here)
* line 3 : date
* line 4 : fractional time
* remaining lines : spectra values
   9.5007600e+01   9.3970000e+01   9.3260500e+01   9.3067300e+01
   8.5000000e+01   8.5000000e+01   8.5000000e+01   8.5000000e+01
   01/07/2012      01/07/2012      01/07/2012      01/07/2012
   2.9940300e+00   3.1491700e+00   3.2526400e+00   3.2804200e+00
   0.0000000e+00   0.0000000e+00   0.0000000e+00   0.0000000e+00
  -1.8073600e+00  -5.0432300e+00  -6.7440700e+00  -4.7440500e+00
  -1.4793200e+00  -3.5119500e+00  -2.1264100e+00  -9.1264000e+00
   ...             ...             ...             ...
\end{Verbatim}
\end{quote}

With column format, a wavelength calibration can be provided with spectra.  It should
be given in first column of each data set.

Other examples of ASCII files supported by QDOAS are given in appendix \textbf{\ref{cha:input-files-format}}.

\section{Installation}
QDOAS packages for Windows, Linux and macOS systems are available in the Conda package manager.  If Conda is not yet installed on your system, you should install it first\footnote{If you are installing Conda specifically for QDOAS, we recommend the lightweight Miniforge distribution available at \url{https://github.com/conda-forge/miniforge}.  Alternatively, you can obtain Miniconda or the full Anaconda distribution from \url{https://anaconda.org}.}
.  Once Conda is installed, run the following command from the Conda prompt to install QDOAS:

\begin{Verbatim}[fontfamily=courier,fontsize=\relsize{-2},xleftmargin=5mm]
conda install conda-forge::qdoas
-\end{Verbatim}

This will install the QDOAS GUI tools qdoas, convolution, ring and usamp, as well as the command line tool \verb|doas_cl|.  Because the GUI tools depend on the Qt libraries, such an installation will take up about 1,4GB of disk space.  Users who only need \verb|doas_cl|  -- e.g.\ for batch processing on a compute server -- can reduce the installation footprint by installing only \verb|doas_cl|:
\begin{Verbatim}[fontfamily=courier,fontsize=\relsize{-2},xleftmargin=5mm]
conda install conda-forge::doas_cl
-\end{Verbatim}

\subsection{Online Help}
QDOAS also provices online help in the form of HTML pages.  The
content is similar to chapters \ref{cha:descr-algor},
\ref{cha:properties} and \ref{cha:tools} of this document.  The HTML
files have to be copied on your disk (usually a subfolder Help of the
directory where executables are installed on Windows systems; a
dedicated subfolder of your system share directory on Linux systems).
The HTML pages are reachable from the Help button of the configuration
dialog boxes.  If the page can not be found, the user is is asked to
provide the location of the help's index.html file, after which the help system should work.

\subsection{Source code}
The source code for QDOAS is freely available at \url{https://github.com/UVVIS-BIRA-IASB/qdoas}.  The README file contains instructions for users who want to compile QDOAS from source.


%%% Local Variables:
%%% mode: latex
%%% TeX-master: "QDOAS_manual"
%%% TeX-PDF-mode: t
%%% End:
